\exercise{9.1}{Three quotients evaluated by the finite limit laws}
\goal Use Theorems~9.2--9.7, with their hypotheses checked, to prove
\begin{align*}
 \text{(a)}\quad &\lim_{n\to\infty}\frac{n+1}{n}=1,\\
 \text{(b)}\quad &\lim_{n\to\infty}\frac{3n+7}{6n-5}=\frac12,\\
 \text{(c)}\quad &\lim_{n\to\infty}
 \frac{17n^5+73n^4-18n^2+3}{23n^5+13n^3}=\frac{17}{23}.
\end{align*}

\idea Divide each polynomial expression by the appropriate positive
power of $n$. The resulting expressions involve constants and
reciprocal powers, whose limits are known. For each use of the quotient
theorem, verify that the denominator terms are nonzero and that their
finite limit is nonzero.

\step{1. Record the finite limit rules being used}
A constant sequence with value $c\in\R$ converges to $c$, because
$\abs{c-c}=0<\varepsilon$ for every $\varepsilon>0$.
Theorem~9.7(a) gives
\[
 \frac1{n^p}\longrightarrow0\qquad(p>0).
\]
The scalar rule is Theorem~9.2, the sum rule is Theorem~9.3,
and the product rule is Theorem~9.4. Subtraction follows by
combining the scalar rule with scalar $-1$ and the sum rule.
The quotient rule, Theorem~9.6, requires a convergent numerator,
a convergent denominator, a nonzero denominator at every index,
and a nonzero denominator limit. All limits used in these rules
are finite real numbers.

\step{2. Part (a): Separate the constant term}
Since $n\in\N$, the denominator $n$ is positive. The identity
\[
 \frac{n+1}{n}=\frac nn+\frac1n=1+\frac1n
\]
therefore holds for every index. The constant sequence $1$ converges
to $1$, and Theorem~9.7(a), with $p=1$, gives $1/n\to0$.
Applying Theorem~9.3 yields
\[
 \frac{n+1}{n}=1+\frac1n\longrightarrow1+0=1.
\]
Here the sum rule applies after rewriting; there is no need to treat
the unbounded sequence $n$ as a finite convergent sequence.

\step{3. Part (b): Normalize numerator and denominator}
For every $n\ge1$, the original denominator satisfies
$6n-5\ge1>0$. Dividing numerator and denominator by $n>0$ gives
\[
 \frac{3n+7}{6n-5}
 =\frac{(3n+7)/n}{(6n-5)/n}
 =\frac{3+7/n}{6-5/n}.
\]
Let $u_n=3+7/n$ and $v_n=6-5/n$.
By Theorem~9.7(a), $1/n\to0$. The scalar rule gives
$7/n\to7\cdot0=0$ and $-5/n\to-5\cdot0=0$.
Adding the constant sequences by Theorem~9.3 gives
\[
 u_n\longrightarrow3+0=3,
 \qquad v_n\longrightarrow6+0=6.
\]
For each $n$, we also have
\[
 v_n=\frac{6n-5}{n}>0.
\]
Thus every denominator term is nonzero, and its limit $6$ is nonzero.
All hypotheses of Theorem~9.6 have now been verified, so
\[
 \frac{3n+7}{6n-5}=\frac{u_n}{v_n}
 \longrightarrow\frac36=\frac12.
\]

\step{4. Part (c): Divide by the largest power in the denominator}
The original denominator is positive for every index:
\[
 23n^5+13n^3=n^3(23n^2+13)>0,
\]
because $n\ge1$. Divide both numerator and denominator by $n^5>0$.
Each term then simplifies as follows:
\[
 \frac{n^4}{n^5}=\frac1n,
 \qquad \frac{n^2}{n^5}=\frac1{n^3},
 \qquad \frac{n^3}{n^5}=\frac1{n^2}.
\]
Consequently
\[
 \frac{17n^5+73n^4-18n^2+3}{23n^5+13n^3}
 =\frac{17+73/n-18/n^3+3/n^5}{23+13/n^2}.
\]
Write the normalized numerator and denominator as $U_n$ and $V_n$.
Theorem~9.7(a), applied with $p=1,2,3,5$, gives
\[
 \frac1n\to0,\qquad \frac1{n^2}\to0,\qquad
 \frac1{n^3}\to0,\qquad \frac1{n^5}\to0.
\]
Multiplication by the indicated constants, using Theorem~9.2, gives
\[
 \frac{73}{n}\to0,\qquad
 -\frac{18}{n^3}\to0,\qquad
 \frac3{n^5}\to0,\qquad
 \frac{13}{n^2}\to0.
\]
Repeated applications of the sum rule, together with the constant
limits, therefore yield
\[
 U_n\longrightarrow17+0+0+0=17,
 \qquad V_n\longrightarrow23+0=23.
\]
Moreover, $V_n=23+13/n^2>23>0$ for every $n$, and its limit
$23$ is nonzero. The quotient rule applies and gives
\[
 \frac{U_n}{V_n}\longrightarrow\frac{17}{23}.
\]

\answer{The three limits are \textup{(a)} $1$,
\textup{(b)} $1/2$, and \textup{(c)} $17/23$.
Every use of the quotient theorem has a denominator that is nonzero
at every index and has a nonzero finite limit.}

\exercise{9.2}{Limits obtained from two given convergent sequences}
\goal Suppose $x_n\to3$, $y_n\to7$, and $y_n\ne0$ for every $n$.
Find
\[
 \text{(a)}\quad\lim_{n\to\infty}(x_n+y_n),
 \qquad
 \text{(b)}\quad\lim_{n\to\infty}\frac{3y_n-x_n}{y_n^2}.
\]

\idea Use the sum rule in part \textup{(a)}. In part \textup{(b)},
build the numerator from scalar multiples and a sum, build the
denominator with the product rule, and then check the quotient rule.

\step{1. Part (a): Add the two limits}
The hypotheses give two finite convergent sequences. Theorem~9.3
therefore applies directly:
\[
 x_n+y_n\longrightarrow3+7=10.
\]

\step{2. Part (b): Evaluate the numerator}
Theorem~9.2 gives
\[
 3y_n\longrightarrow3\cdot7=21,
 \qquad -x_n\longrightarrow(-1)\cdot3=-3.
\]
Theorem~9.3 then gives
\[
 3y_n-x_n=3y_n+(-x_n)\longrightarrow21+(-3)=18.
\]
The numerator here is $3y_n-x_n$; only the denominator contains a
square.

\step{3. Part (b): Evaluate the denominator and take the quotient}
Since $y_n\to7$, Theorem~9.4 applied to the two factors $y_n$ gives
\[
 y_n^2=y_n\,y_n\longrightarrow7\cdot7=49.
\]
The assumption $y_n\ne0$ makes $y_n^2>0$ at every index.
Thus every denominator is nonzero, and its limit $49$ is also
nonzero. Theorem~9.6 now gives
\[
 \frac{3y_n-x_n}{y_n^2}\longrightarrow\frac{18}{49}.
\]

\answer{The limits are \textup{(a)} $10$ and
\textup{(b)} $18/49$.}

\exercise{9.3}{A polynomial expression in convergent sequences}
\goal Assume $a_n\to a\in\R$ and $b_n\to b\in\R$. For
\[
 s_n=\frac{a_n^3+4a_n}{b_n^2+1},
\]
prove carefully that
\[
 s_n\longrightarrow\frac{a^3+4a}{b^2+1}.
\]

\idea Form the powers one multiplication at a time. Then use the
scalar and sum rules. The added $1$ makes the denominator and its
limit positive even when $b_n=0$ or $b=0$.

\step{1. Construct the square and cube of the numerator sequence}
The two factors in $a_n\,a_n$ both converge to $a$.
The product rule, Theorem~9.4, gives
\[
 a_n^2=a_n\,a_n\longrightarrow a\cdot a=a^2.
\]
We may now apply the product rule again, this time to $a_n^2$ and
$a_n$, since both have finite limits:
\[
 a_n^3=(a_n^2)a_n\longrightarrow a^2\cdot a=a^3.
\]
Separately, Theorem~9.2 gives $4a_n\to4a$.
The sum rule therefore yields
\[
 a_n^3+4a_n\longrightarrow a^3+4a.
\]

\step{2. Construct the denominator limit}
Apply Theorem~9.4 to $b_n\,b_n$ to obtain
\[
 b_n^2\longrightarrow b\cdot b=b^2.
\]
The constant sequence $1$ converges to $1$, so Theorem~9.3 gives
\[
 b_n^2+1\longrightarrow b^2+1.
\]
For every real $b_n$, its square is nonnegative, and thus
$b_n^2+1\ge1>0$. The denominator terms are consequently nonzero.
Likewise $b^2+1\ge1>0$, so their limit is nonzero.

\step{3. Apply the quotient theorem with all hypotheses established}
Step~1 proves convergence of the numerator, Step~2 proves convergence
of the denominator, and the positivity checks verify both nonzero
requirements. Theorem~9.6 therefore gives the desired limit.

\answer{$\displaystyle
 \lim_{n\to\infty}s_n=\frac{a^3+4a}{b^2+1}$
for all finite real limits $a$ and $b$, including $b=0$.}

\exercise{9.4}{A nested-square-root recursion}
\goal For $s_1=1$ and $s_{n+1}=\sqrt{s_n+1}$, list the first
four terms and, assuming convergence, prove that the limit is
$(1+\sqrt5)/2$.

\idea First check that every radical is defined and every term is
at least $1$. Under the stated convergence assumption, square the
recursion and apply the product and sum limit laws. The lower bound
will select the appropriate quadratic root.

\step{1. Verify that the recursion stays defined and positive}
The initial value is $s_1=1\ge1$. If $s_n\ge1$, then
$s_n+1\ge2>0$, so its nonnegative real square root exists.
Moreover, $\sqrt{s_n+1}\ge1$: a nonnegative number less than $1$
would have square less than $1$, whereas this square is at least $2$.
Thus $s_{n+1}\ge1$. Induction shows that all terms are defined
and satisfy $s_n\ge1$.

\step{2. Part (a): Substitute successively four times}
The successive terms are
\begin{align*}
 s_1&=1,\\
 s_2&=\sqrt{s_1+1}=\sqrt{1+1}=\sqrt2,\\
 s_3&=\sqrt{s_2+1}=\sqrt{\sqrt2+1}
      =\sqrt{1+\sqrt2},\\
 s_4&=\sqrt{s_3+1}=\sqrt{\sqrt{1+\sqrt2}+1}
      =\sqrt{1+\sqrt{1+\sqrt2}}.
\end{align*}

\step{3. A shift of a convergent sequence has the same limit}
We will use this elementary observation for this and the next two
recursions. If $u_n\to u\in\R$, then $u_{n+1}\to u$.
Indeed, given $\varepsilon>0$, choose $N$ so that
$m>N$ implies $\abs{u_m-u}<\varepsilon$.
If $n>N$, then also $n+1>N$, so
$\abs{u_{n+1}-u}<\varepsilon$. This is exactly the limit
definition for the shifted sequence.

\step{4. Part (b): Record the restriction on a finite limit}
Assume, as requested, that $s_n\to L\in\R$.
Since $s_n\ge1$ for every $n$, we have $L\ge1$.
To see the bound directly, if $L<1$, set
$\varepsilon=(1-L)/2>0$. Eventually
$\abs{s_n-L}<\varepsilon$, giving
\[
 s_n<L+\varepsilon=\frac{L+1}{2}<1,
\]
contrary to the bound on the terms. This is also the eventual-bound
principle of Exercise~8.9(a).

\step{5. Pass to limits in a squared identity}
The defining radical equality implies
\[
 s_{n+1}^2=s_n+1
\]
for every $n$. By Step~3, $s_{n+1}\to L$.
Theorem~9.4 applied to $s_{n+1}\,s_{n+1}$ gives
$s_{n+1}^2\to L^2$. Theorem~9.3 gives
$s_n+1\to L+1$. These two expressions are equal term by term,
so uniqueness of finite limits, proved at the end of Section~7,
gives
\[
 L^2=L+1.
\]
Only the product and sum laws were used; a square-root limit rule
is not needed.

\step{6. Solve the quadratic equation and select its permitted root}
Rearrange and complete the square:
\[
 L^2-L-1=0,
 \qquad 4L^2-4L=4,
 \qquad (2L-1)^2=5.
\]
Consequently $2L-1=\sqrt5$ or $2L-1=-\sqrt5$, which gives
\[
 L=\frac{1+\sqrt5}{2}
 \quad\text{or}\quad L=\frac{1-\sqrt5}{2}.
\]
We have $\sqrt5>1$, since the opposite nonnegative inequality
would imply $5\le1$ after squaring. Thus
$(1-\sqrt5)/2<0$, contradicting $L\ge1$.
The only remaining root is $(1+\sqrt5)/2$.

\answer{The first four terms are
$1$, $\sqrt2$, $\sqrt{1+\sqrt2}$, and
$\sqrt{1+\sqrt{1+\sqrt2}}$.
Under the given convergence assumption,
$\displaystyle\lim s_n=(1+\sqrt5)/2$.}

\exercise{9.5}{A positive rational recursion}
\goal Let $t_1=1$ and
\[
 t_{n+1}=\frac{t_n^2+2}{2t_n}\qquad(n\ge1).
\]
Assuming that $(t_n)$ converges to a finite real number, find its limit.

\idea Prove positivity so that every term is defined, and strengthen
it to the uniform lower bound $t_n\ge1$. The lower bound makes the
assumed limit nonzero before the quotient theorem is applied.

\step{1. Check that the recursion is defined at every index}
We start with $t_1=1>0$. If $t_n>0$, then its denominator
$2t_n$ is positive and its numerator $t_n^2+2$ is positive.
The quotient is therefore a well-defined positive real number.
Induction proves $t_n>0$ for all $n$, and in particular no
denominator in the recursion vanishes.

\step{2. Obtain a positive uniform lower bound}
For every $n$, use the positive denominator to compute
\begin{align*}
 t_{n+1}-1
 &=\frac{t_n^2+2}{2t_n}-\frac{2t_n}{2t_n}\\
 &=\frac{t_n^2-2t_n+2}{2t_n}\\
 &=\frac{(t_n-1)^2+1}{2t_n}>0.
\end{align*}
The last numerator is positive because $(t_n-1)^2\ge0$.
Thus every term after $t_1$ is greater than $1$, and $t_1=1$.
It follows that $t_n\ge1$ for every $n$.

\step{3. The convergence assumption gives a nonzero limit}
Suppose $t_n\to L\in\R$. The bound $t_n\ge1$ passes to the
limit by Exercise~8.9(a), so $L\ge1>0$.
Equivalently, if $L<1$, the positive tolerance $(1-L)/2$
would force eventual terms to be less than
$L+(1-L)/2=(L+1)/2<1$, a contradiction.
In particular, $2L\ne0$.

\step{4. Apply the finite limit laws to the recursion}
The product rule, Theorem~9.4, gives
\[
 t_n^2=t_n\,t_n\longrightarrow L^2.
\]
The constant sequence $2$ converges to $2$, and hence the sum rule
gives $t_n^2+2\to L^2+2$. The scalar rule gives
$2t_n\to2L$. The terms $2t_n$ are nonzero by Step~1,
and their limit $2L$ is nonzero by Step~3. Thus Theorem~9.6 applies:
\[
 \frac{t_n^2+2}{2t_n}\longrightarrow\frac{L^2+2}{2L}.
\]
The shift observation in Exercise~9.4 also gives $t_{n+1}\to L$.
The recursion identifies these two sequences, so uniqueness of
finite limits gives
\[
 L=\frac{L^2+2}{2L}.
\]

\step{5. Solve and select the positive root}
Since $2L\ne0$, multiplying the last equality by $2L$ is valid:
\[
 2L^2=L^2+2,
 \qquad L^2=2.
\]
Thus $L=\sqrt2$ or $L=-\sqrt2$.
The root $-\sqrt2$ is negative and violates $L\ge1$.
Consequently the assumed finite limit is $\sqrt2$.

\answer{Under the convergence assumption,
$\displaystyle\lim_{n\to\infty}t_n=\sqrt2$.
The nonzero denominator limit was established before the quotient
theorem was used.}

\exercise{9.6}{A necessary finite-limit equation and an actual infinite limit}
\goal For $x_1=1$ and $x_{n+1}=3x_n^2$, address all three questions:
\textup{(a)} determine the possibilities for a finite real limit;
\textup{(b)} decide whether the actual sequence has a limit;
\textup{(c)} explain why the conclusions are consistent.

\idea A hypothetical finite limit must satisfy the recursion's
fixed-point equation. Separately, use an induction estimate for the
actual terms. Distinguish finite convergence from the extended
infinite-limit convention in Definition~9.8.

\step{1. Part (a): Assume a finite limit before applying the product rule}
Suppose $x_n\to a\in\R$. The shift observation proved in
Exercise~9.4 gives $x_{n+1}\to a$.
Theorem~9.4 gives $x_n^2\to a^2$, and Theorem~9.2 then gives
$3x_n^2\to3a^2$. Since $x_{n+1}=3x_n^2$ for every $n$,
uniqueness of finite limits implies
\[
 a=3a^2.
\]
Subtracting $a$ and factoring gives
\[
 0=3a^2-a=a(3a-1).
\]
The zero-product property in $\R$ says that one of the two factors
must be zero. Thus $a=0$ or $3a-1=0$; the second equation gives
$a=1/3$ after division by $3\ne0$.

This proves only a necessary condition: \emph{if the sequence has a
finite real limit, that limit must be $0$ or $1/3$}. The computation
does not prove that such a finite limit exists.

\step{2. Part (b): Bound the actual terms below by their indices}
We prove $x_n\ge n$ for every $n\in\N$.
For $n=1$, the equality $x_1=1$ supplies the base case.
Suppose $x_n\ge n$ for some $n\ge1$.
In particular $x_n\ge1$, so
$x_n^2-x_n=x_n(x_n-1)\ge0$. It follows that
\[
 x_{n+1}=3x_n^2\ge3x_n\ge3n\ge n+1.
\]
The final inequality holds because
$3n-(n+1)=2n-1\ge1>0$ for $n\ge1$.
This completes the induction.

Now let $M>0$ and take the threshold $N=M$.
For every integer $n>N$, the bound gives
\[
 x_n\ge n>M.
\]
This is exactly Definition~9.8 for $x_n\to+\infty$.
Thus the actual sequence has the extended limit $+\infty$.

It cannot have a finite real limit. For a direct check, if
$x_n\to a\in\R$, then eventually $\abs{x_n-a}<1$,
which implies $x_n<a+1$. Choose an integer $n$ larger than
both that eventual threshold and $a+1$, using the Archimedean
property. The bound $x_n\ge n$ would instead give $x_n>a+1$,
a contradiction. This is also consistent with Theorem~9.1,
which says every finitely convergent sequence is bounded.

\step{3. Part (c): Explain the two meanings of the limit statement}
Part \textup{(a)} began with the hypothesis of a \emph{finite real}
limit $a$. Its use of the finite product and scalar limit laws is
valid under that hypothesis. Part \textup{(b)} shows that the
hypothesis fails for this initial value: the terms grow without bound.
Consequently the necessary alternatives in part \textup{(a)} are
not attained by this sequence.

Under the convention introduced after Definition~9.8, saying that
``the limit exists'' also allows divergence to $+\infty$ or
$-\infty$. In that stated convention the answer to part \textup{(b)}
is yes, with limit $+\infty$; the sequence does not converge to a
finite real number. Theorems~9.2--9.6 were proved for finite real
limits, so their product calculation in part \textup{(a)} cannot be
applied with $a=+\infty$. The extended infinite-limit conclusion
therefore supplies no contradiction to the finite necessary condition.

\answer{\textup{(a)} Any hypothetical finite real limit is $0$ or
$1/3$. \textup{(b)} The actual sequence satisfies
$\displaystyle\lim x_n=+\infty$ in the sense of Definition~9.8,
and has no finite real limit. \textup{(c)} A necessary condition
for finite convergence does not assert its existence, and finite
product laws cannot be applied at an infinite limit.}

\exercise{9.7}{Finish the estimate for the root of the index}
\goal Complete the argument that
\[
 s_n=n^{1/n}-1\longrightarrow0,
\]
and hence finish the proof of Theorem~9.7(c), $n^{1/n}\to1$.

\idea Use the nonnegative binomial terms in $(1+s_n)^n$ to obtain
an explicit bound on $s_n$. Then choose a threshold directly in
terms of the prescribed positive tolerance.

\step{1. Verify the sign and the defining power identity}
Let $u_n=n^{1/n}$ denote the positive real $n$th root of $n$.
Then $u_n^n=n\ge1$. If $u_n<1$, the positivity $0<u_n<1$
would give $u_n^n<1$, a contradiction. Hence $u_n\ge1$ and
\[
 s_n=u_n-1\ge0,
 \qquad (1+s_n)^n=u_n^n=n.
\]
For $n=1$, this gives $s_1=1-1=0$.

\step{2. Derive the binomial bound for $n\ge2$}
For an integer $n\ge2$, the binomial theorem gives
\[
 n=(1+s_n)^n
   =1+ns_n+\frac{n(n-1)}2s_n^2
      +\sum_{j=3}^{n}\binom nj s_n^j.
\]
When $n=2$, the final sum is empty and equals $0$.
All terms of that sum are nonnegative, since $s_n\ge0$ and all
the binomial coefficients are nonnegative. Also $ns_n\ge0$,
and the constant term $1$ is strictly positive. Therefore
\[
 n>\frac{n(n-1)}2s_n^2.
\]
Dividing by the positive number $n$ gives
$1>(n-1)s_n^2/2$. Multiplying by $2/(n-1)>0$ yields
\[
 s_n^2<\frac2{n-1}.
\]
Both $s_n$ and the square root of the right-hand side are
nonnegative, so the order of their squares gives
\[
 0\le s_n<\sqrt{\frac2{n-1}}\qquad(n\ge2).
\]

\step{3. Specify the threshold and verify the limit definition}
Let $\varepsilon>0$ be arbitrary and set
\[
 N=\max\left\{2,\ 1+\frac2{\varepsilon^2}\right\}.
\]
A real threshold is permitted by the definition of convergence;
an integer larger than it may also be chosen by the Archimedean
property. If $n>N$, then $n\ge2$ and
\[
 n-1>\frac2{\varepsilon^2}>0.
\]
Taking reciprocals of positive quantities and multiplying by $2$
gives $2/(n-1)<\varepsilon^2$. As $\varepsilon>0$, this implies
\[
 \sqrt{\frac2{n-1}}<\varepsilon.
\]
Using the bound from Step~2, we conclude
\[
 \abs{s_n-0}=s_n<\sqrt{\frac2{n-1}}<\varepsilon
 \qquad(n>N).
\]
This is the required proof that $s_n\to0$.

\step{4. Return to the root sequence}
The constant sequence $1$ converges to $1$, and $s_n\to0$ was
proved without assuming a limit for $n^{1/n}$.
The sum rule, Theorem~9.3, therefore gives
\[
 n^{1/n}=1+s_n\longrightarrow1+0=1.
\]

\answer{The explicit estimate
$0\le n^{1/n}-1<\sqrt{2/(n-1)}$ for $n\ge2$, together with
$N=\max\{2,1+2/\varepsilon^2\}$, proves
$n^{1/n}-1\to0$ and hence $n^{1/n}\to1$.}
